% GENERATED FILE. Edit canonical lessons, then run npm run generate:books.
% canonical-source-hash: fnv1a64:e09d86913e18e401
% canonical-lessons: GU-C412-dero, GU-C412-padav, GU-C412-ratvaso, GU-C412-khodvu-peg, GU-C412-galvu-spend

\chapter{Tents and Camps}
\label{ch:gu-tents-and-camps}
\hlchaptermodality{412}

\begin{chapteropening}
\textbf{By the end of this chapter:} \emph{I can talk about a night spent in a camp.}

Complete the last lesson of chapter 412: I can talk about a night spent in a camp.
\end{chapteropening}

\section[\texorpdfstring{\gu{ડેરો} (ḍero)}{ડેરો (ḍero)}]{\gu{ડેરો} (ḍero) — a tent; a camp (for a stay)}

\label{lesson:GU-C412-dero}



\begin{quote}
Before the new one: say the Gujarati for to roll up, then the Gujarati for to put to sleep.
\end{quote}

\subsection*{You'll want to know: \gu{ડેરો}}
\textbf{\gu{ડેરો}} — \emph{ḍero} — \textquotedblleft{}a tent; a camp (for a stay)\textquotedblright{}.

\begin{quote}
\gu{નદીના} \gu{કિનારે} \gu{અમે} \gu{ડેરો} \gu{નાખ્યો}.
\end{quote}

\begin{quote}
\emph{nadīnā kināre ame ḍero nākhyo.}
\end{quote}

\begin{quote}
We pitched our tent on the river bank.
\end{quote}

A masculine noun.

\begin{cousinweb}
\textbf{\gu{ડેરો}} (\emph{ḍero}) goes back to \emph{*ḍēra}, \textquotedblleft{}a resting place\textquotedblright{}, a form CDIAL reconstructs; CDIAL lists the Gujarati word, \textquotedblleft{}a tent, a shelter\textquotedblright{} (CDIAL 5564).
\end{cousinweb}

\begin{grammarlens}[title={What you've built}]
One more word for getting around.
\end{grammarlens}

\subsection*{Your turn}
\begin{itemize}
\raggedright
  \item \emph{Say it:} \emph{ḍero}
  \item \emph{Say it:} \emph{ḍero}, once more
  \item \emph{Say it:} \emph{nadīnā kināre ame ḍero nākhyo.} — We pitched our tent on the river bank.
  \item \emph{Recall:} say the Gujarati for to roll up, then the Gujarati for to put to sleep, then say \emph{ḍero} again
\end{itemize}

\begin{tcolorbox}[breakable,colback=teal!4,colframe=teal!35!black,title={Before you move on}]
What does \gu{ડેરો} mean? (\textquotedblleft{}A tent; a camp (for a stay)\textquotedblright{}.) Say it once more.
\end{tcolorbox}
\section[\texorpdfstring{\gu{પડાવ} (paḍāv)}{પડાવ (paḍāv)}]{\gu{પડાવ} (paḍāv) — a halt for the night; a camp}

\label{lesson:GU-C412-padav}



\begin{quote}
Before the new one: say the Gujarati for to put to sleep, then the Gujarati for a tent.
\end{quote}

\subsection*{You'll want to know: \gu{પડાવ}}
\textbf{\gu{પડાવ}} — \emph{paḍāv} — \textquotedblleft{}a halt for the night; a camp\textquotedblright{}.

\begin{quote}
\gu{સાંજે} \gu{અમારો} \gu{પડાવ} \gu{ગામની} \gu{બહાર} \gu{હતો}.
\end{quote}

\begin{quote}
\emph{sā̃je amāro paḍāv gāmnī bahār hato.}
\end{quote}

\begin{quote}
In the evening our camp was outside the village.
\end{quote}

A masculine noun.

\begin{cousinweb}
\textbf{\gu{પડાવ}} (\emph{paḍāv}) seems to be built on \gu{પડવું}, to fall, here to stop and lie down for the night; CDIAL lists that verb under Sanskrit \emph{pátati}, \textquotedblleft{}falls\textquotedblright{} (CDIAL 7722), but gives no form for \gu{પડાવ} itself, so the link is hedged.
\end{cousinweb}

\begin{grammarlens}[title={What you've built}]
One more word for getting around.
\end{grammarlens}

\subsection*{Your turn}
\begin{itemize}
\raggedright
  \item \emph{Say it:} \emph{paḍāv}
  \item \emph{Say it:} \emph{paḍāv}, once more
  \item \emph{Say it:} \emph{sā̃je amāro paḍāv gāmnī bahār hato.} — In the evening our camp was outside the village.
  \item \emph{Recall:} say the Gujarati for to put to sleep, then the Gujarati for a tent, then say \emph{paḍāv} again
\end{itemize}

\begin{tcolorbox}[breakable,colback=teal!4,colframe=teal!35!black,title={Before you move on}]
What does \gu{પડાવ} mean? (\textquotedblleft{}A halt for the night; a camp\textquotedblright{}.) Say it once more.
\end{tcolorbox}
\section[\texorpdfstring{\gu{રાતવાસો} (rātvāso)}{રાતવાસો (rātvāso)}]{\gu{રાતવાસો} (rātvāso) — a night's stay}

\label{lesson:GU-C412-ratvaso}



\begin{quote}
Before the new one: say the Gujarati for a tent, then the Gujarati for a halt for the night.
\end{quote}

\subsection*{You'll want to know: \gu{રાતવાસો}}
\textbf{\gu{રાતવાસો}} — \emph{rātvāso} — \textquotedblleft{}a night's stay\textquotedblright{}.

\begin{quote}
\gu{અમે} \gu{ધર્મશાળામાં} \gu{રાતવાસો} \gu{કર્યો}.
\end{quote}

\begin{quote}
\emph{ame dharmaśāḷāmā̃ rātvāso karyo.}
\end{quote}

\begin{quote}
We stayed the night at the rest house.
\end{quote}

A masculine noun.

\begin{cousinweb}
\textbf{\gu{રાતવાસો}} (\emph{rātvāso}) is built from \gu{રાત}, night, and \gu{વાસો}, a stay, from Sanskrit \emph{vāsa}, \textquotedblleft{}a dwelling, an abode\textquotedblright{} (Monier-Williams): a roof for one night.
\end{cousinweb}

\begin{grammarlens}[title={What you've built}]
One more word for getting around.
\end{grammarlens}

\subsection*{Your turn}
\begin{itemize}
\raggedright
  \item \emph{Say it:} \emph{rātvāso}
  \item \emph{Say it:} \emph{rātvāso}, once more
  \item \emph{Say it:} \emph{ame dharmaśāḷāmā̃ rātvāso karyo.} — We stayed the night at the rest house.
  \item \emph{Recall:} say the Gujarati for a tent, then the Gujarati for a halt for the night, then say \emph{rātvāso} again
\end{itemize}

\begin{tcolorbox}[breakable,colback=teal!4,colframe=teal!35!black,title={Before you move on}]
What does \gu{રાતવાસો} mean? (\textquotedblleft{}A night's stay\textquotedblright{}.) Say it once more.
\end{tcolorbox}
\section[\texorpdfstring{\gu{ખોડવું} (khoḍvũ)}{ખોડવું (khoḍvũ)}]{\gu{ખોડવું} (khoḍvũ) — to drive in (a peg); to plant (a pole)}

\label{lesson:GU-C412-khodvu-peg}



\begin{quote}
Before the new one: say the Gujarati for a halt for the night, then the Gujarati for a night's stay.
\end{quote}

\subsection*{You'll want to know: \gu{ખોડવું}}
\textbf{\gu{ખોડવું}} — \emph{khoḍvũ} — \textquotedblleft{}to drive in (a peg); to plant (a pole)\textquotedblright{}.

\begin{quote}
\gu{છોકરાએ} \gu{રેતીમાં} \gu{લાંબી} \gu{ડાળી} \gu{ખોડી}.
\end{quote}

\begin{quote}
\emph{chokrāe retīmā̃ lā̃bī ḍāḷī khoḍī.}
\end{quote}

\begin{quote}
The boy stuck a long branch into the sand.
\end{quote}

Past, with \gu{તેણે} (he or she …): \gu{તેણે} \gu{ડાળી} \gu{ખોડી}, \emph{teṇe ḍāḷī khoḍī}.

\begin{cousinweb}
\textbf{\gu{ખોડવું}} (\emph{khoḍvũ}) may belong with a Gujarati word for a large block of wood, which CDIAL lists under Sanskrit \emph{kṣōḍa}, \textquotedblleft{}a post to which an elephant is tied\textquotedblright{} (CDIAL 3748), beside Marathi \emph{khoḍ}, \textquotedblleft{}a tree stump\textquotedblright{}: to drive a post into the ground. CDIAL gives no form for \gu{ખોડવું} itself, so the link is hedged.
\end{cousinweb}

\begin{grammarlens}[title={What you've built}]
One more word for getting around.
\end{grammarlens}

\subsection*{Your turn}
\begin{itemize}
\raggedright
  \item \emph{Say it:} \emph{khoḍvũ}
  \item \emph{Say it:} \emph{khoḍvũ}, once more
  \item \emph{Say it:} \emph{chokrāe retīmā̃ lā̃bī ḍāḷī khoḍī.} — The boy stuck a long branch into the sand.
  \item \emph{Recall:} say the Gujarati for a halt for the night, then the Gujarati for a night's stay, then say \emph{khoḍvũ} again
\end{itemize}

\begin{tcolorbox}[breakable,colback=teal!4,colframe=teal!35!black,title={Before you move on}]
What does \gu{ખોડવું} mean? (\textquotedblleft{}To drive in (a peg); to plant (a pole)\textquotedblright{}.) Say it once more.
\end{tcolorbox}
\section[\texorpdfstring{\gu{ગાળવું} (gāḷvũ)}{ગાળવું (gāḷvũ)}]{\gu{ગાળવું} (gāḷvũ) — to spend (time)}

\label{lesson:GU-C412-galvu-spend}



\begin{quote}
Before the new one: say the Gujarati for a night's stay, then the Gujarati for to drive in.
\end{quote}

\subsection*{You'll want to know: \gu{ગાળવું}}
\textbf{\gu{ગાળવું}} — \emph{gāḷvũ} — \textquotedblleft{}to spend (time)\textquotedblright{}.

\begin{quote}
\gu{અમે} \gu{બે} \gu{દિવસ} \gu{ગામમાં} \gu{ગાળ્યા}.
\end{quote}

\begin{quote}
\emph{ame be divas gāmmā̃ gāḷyā.}
\end{quote}

\begin{quote}
We spent two days in the village.
\end{quote}

Past, with \gu{તેણે} (he or she …): \gu{તેણે} \gu{રાત} \gu{ત્યાં} \gu{ગાળી}, \emph{teṇe rāt tyā̃ gāḷī}.

\begin{cousinweb}
\textbf{\gu{ગાળવું}} (\emph{gāḷvũ}) is a tadbhava of Sanskrit \emph{gālayati}, \textquotedblleft{}causes to drop, strains, melts\textquotedblright{}; CDIAL lists the Gujarati verb, which also means to pass time (CDIAL 4144): the days drip away.
\end{cousinweb}

\begin{grammarlens}[title={What you've built}]
One more word for getting around.
\end{grammarlens}

\subsection*{Your turn}
\begin{itemize}
\raggedright
  \item \emph{Say it:} \emph{gāḷvũ}
  \item \emph{Say it:} \emph{gāḷvũ}, once more
  \item \emph{Say it:} \emph{ame be divas gāmmā̃ gāḷyā.} — We spent two days in the village.
  \item \emph{Recall:} say the Gujarati for a night's stay, then the Gujarati for to drive in, then say \emph{gāḷvũ} again
\end{itemize}

\begin{tcolorbox}[breakable,colback=teal!4,colframe=teal!35!black,title={Before you move on}]
What does \gu{ગાળવું} mean? (\textquotedblleft{}To spend (time)\textquotedblright{}.) Say it once more.
\end{tcolorbox}
